\head{EGA V 短篇：函数的奇异零点与超奇异零点}
\label{ega:V-short}
\source{\texttt{EGAV1-2.ps}，转换后共6页：标题页及正文印刷第1--5页；与\texttt{EGAVaiello.pdf}第62--66页前部同文}
\noindent\textbf{本版同文校订。}定义1至注4(c)合入Dot随后两批的校订；重合段落只排印一次，原有注4(d)及后文完整保留。两载体的接合位置与仍待核的数学问题见本版核查报告。
\noindent\textbf{载体标题。} EGA V：§1及§§2.15、2.16（原编号为EGA IV：§16及§§17.15、17.16）。Grothendieck 的 EGA V“预备笔记”的编辑英译，由 Piotr Blass 与 Joseph Blass 完成。

\egasection{V.1（原IV.16）：函数的奇异集、超奇异集与微分判据}
这一节将在关于超平面截面的V.5（原IV.20）中使用，不过在我看来，它自然的位置是在这里。

\egas{1}{定义}
设$X$为正则预概形，$\varphi$为$\mathcal O_X$的一个截面。若点$x\in X$满足
\[
\varphi_x\in\mathfrak m_x^2,
\]
则称$x$为$\varphi$的一个\emph{奇异零点（或根）}。\cmt{V-H-C16-01}{底本写作$\mathfrak m_X$，但所讨论的是点$x$处的函数芽及余域$k(x)$；故明确记$\mathfrak m_x$为局部环$\mathcal O_{X,x}$的极大理想，$\varphi_x$为截面在该局部环中的芽。}
若$x$是一个奇异零点（或根），并且$\varphi_x$所确定的元素
\[
[\varphi_x]\in\mathfrak m_x^2/\mathfrak m_x^3
\simeq\operatorname{Sym}^2_{k(x)}(\mathfrak m_x/\mathfrak m_x^2)
\]
在被解释为$k(x)$上向量空间$\mathfrak m_x/\mathfrak m_x^2$的对偶空间$t_x$上的二次型时，是一个退化二次型，则称$x$为一个\emph{超奇异零点}。\cmt{V-H-C16-02}{底本右端实际印为$\operatorname{Sym}(\mathfrak m_X/\mathfrak m_X^2)$，未写二次分次。左端是极大理想滤过的二次分次，且下文明确解释为二次型，所以此处补为$\operatorname{Sym}^2$；所用同构是正则局部环的伴随分次代数同构的二次部分。关于特征2，须结合合并英文稿第73--74页的原作者术语修订及第78页（16.1）--（16.2）：后文以相关对称双线性型的判别式规定“非退化”，并把它与“光滑二次型”区别。这里的超奇异术语不能据通常的普通二次奇点用法自行替换。后续定义体系尚未连续译入本短篇。}
（不是超奇异零点的奇异零点〔或根〕，称为\emph{奇异而非超奇异的零点}。）\cmt{D17-S13}{本处原稿称“普通奇异零点”。原作者在合并英文稿第73--74页明确撤回这种命名，改称“奇异而非超奇异的零点”或“非退化二次奇异零点”。特征2下$X^2+YZ$是作者举出的反例：它给出通常意义的普通二次奇点，而相关三变量交替双线性型必退化。本版只在此落实已经明确的术语订正，未把后续第17节及附录的全文计作已译。}

\egas{2}{注}
若$x\in V(\varphi)$，则$x$是$\varphi$的一个非奇异零点，当且仅当$\varphi_x\ne0$且$x$是一个非奇异点，也就是说，$x$是$V(\varphi)$的一个正则点；换言之，当且仅当$x$是$V(\varphi)$的正则点，并且$V(\varphi)$在$x$处的局部子概形不等于$X$在$x$处的局部概形。\cmt{V-H-C16-03}{原句中“a non-singular point”由紧随的“i.e.”解释为$V(\varphi)$的正则点，而不是仅指$X$的正则点。此处仍在定义1的正则环境中；$V(\varphi)$按由$\varphi$定义的闭子概形理解。条件$\varphi_x\ne0$指局部环中的芽不为零，不是余域中的值不为零：既然$x\in V(\varphi)$，其余域值恰为零。原文末尾写“在$x$的邻域内$V(\varphi)\ne X$”；若误解为存在一个邻域使两者不等，便不能推出芽非零。本版明确写成局部子概形不等，保留原文解释链中的必要非零条件；反例及量词说明见本版报告D17-S-C03。}

\egas{3}{定义}
设$X$为域$k$上的光滑预概形，$\varphi$为$\mathcal O_X$的一个截面，且$x\in V(\varphi)$。如果对于$k$的每一个扩域$k'$，以及$X$的每一个以$k'$为值域、位于$x$的点$\xi$，在$X_{k'}$中对应的点$x'$都是$\varphi_{k'}$的奇异零点（分别为超奇异零点），则称$x$为$\varphi$相对于$k$的一个\emph{几何奇异零点}（分别为\emph{几何超奇异零点}）。\cmt{V-H-C16-04}{底本在“for every extension $k'$ of $k$”之后印作“with values in $k$”，与所述基变换及注4(a)不合；此处校为“以$k'$为值域”。其含义是一个$k$态射$\xi:\operatorname{Spec}k'\to X$，像为$x$；它诱导$X_{k'}=X\times_k\operatorname{Spec}k'$中的点$x'$。并非任意扩域都存在这样的$\xi$；量词只要求检验实际存在的点。底本写$X'_k$和$\varphi'_k$，这里统一以$X_{k'}$、$\varphi_{k'}$表示基变换及拉回截面，避免撇号位置混淆。}

\egas{4}{注}
\begin{enumerate}
\item[(a)] 由后面将建立的判据可知，在定义3中，只须用一个以某个$k'$为值域的点作检验即可。例如，可以取$k'=k(x)$或$\overline{k(x)}$，并取以这样的$k'$为值域的典范点。\cmt{V-H-C16-05}{底本第二个$k(x)$上有横线，此处明确为余域的代数闭包$\overline{k(x)}$。典范点来自$\operatorname{Spec}k(x)\to X$，或其沿$k(x)\hookrightarrow\overline{k(x)}$的复合。作者以尚待展开的判据为依据；本段译校未补造该判据及其证明。}
\item[(b)] 由注2可知，$x$是$\varphi$的几何非奇异零点，当且仅当$\varphi_x\ne0$，且$V(\varphi)$在$x$处光滑。\cmt{V-H-C16-06}{这里仍在定义3的域$k$上讨论，因此“光滑”指在$k$上光滑。$\varphi_x\ne0$是局部环中的芽非零，并不表示在余域中的函数值非零。$x\in V(\varphi)$时后者为零；该非零条件不能省略，否则零截面也会被误判为非奇异。}
\item[(c)] 假设预概形$X$在另一个预概形$Y$上光滑，$\varphi$为$\mathcal O_X$的一个截面，且$x\in V(\varphi)$。若$x$在纤维$X_s$上是相对于$k(s)$的奇异零点（或超奇异零点），就称$x$为相对于$Y$的奇异零点（或超奇异零点）；这里$s$是$x$在$Y$中的像。\cmt{V-H-C16-07}{纤维上所用函数是$\varphi$的限制。底本此项写“singular zero (or supersingular) relative to $Y$”和“such relative to $k(s)$”，没有在这两处重复“geometrically”；正文保留这一措辞。在紧接定义3的语境中，“相对于$k(s)$”承接域上几何版本的定义，不应悄然改成仅检验纤维局部环的未作扩域的版本。术语是否在后文统一，仍须继续对照。}
\item[(d)] 在定义1的条件下，若将 $\varphi$ 换成 $\varphi'=u\varphi$，其中 $u$ 在 $x$ 处可逆，则 $x\in V(\varphi)$ 的奇异性或超奇异性不变。因而定义1以及定义3都可以自然地推广到 $\varphi$ 为某个可逆模 $L$ 的截面的情形，并在 $L=\OO_X$ 时恢复原定义。
\end{enumerate}

设预概形 $X$ 对预概形 $Y$ 光滑，$\varphi$ 为 $\OO_X$ 的截面。它给出 $\PP^2_{X/Y}$ 的截面 $d^2_{X/Y}\varphi$，后者约化为 $\PP^1_{X/Y}$ 的截面 $d^1_{X/Y}\varphi$，再约化为 $\PP^0_{X/Y}=\OO_X$ 的截面 $d^0\varphi=\varphi$。

\egas{5}{命题}
$d^0\varphi$ 的零点集是 $\varphi$ 的零点集 $V(\varphi)$；$d^1\varphi$ 的零点集是 $\varphi$ 相对于 $Y$ 的奇异零点集 $V(\varphi)_{\mathrm{sing}}$。
\cmt{D17-S01}{底本此句写相对于$S$，但前提中的相对底为$Y$，且$S$尚未引入；本版据前提订正。原历史注号为T08。}
\begin{proof}[原文证明]
第一断言是平凡的。第二断言就是 Jacobian 判据；也可由如下典范同构得出：若 $x$ 是 $k$-预概形 $X$ 的一个 $k$-有理点，则
\[
 \mm_x/\mm_x^2\simeq\Omega^1_{X/k}(x).
\]
\end{proof}
注意 $\operatorname{gr}^1(\PP^1_{X/Y})\simeq\Omega^1_{X/Y}$，所以限制 $d^1\varphi|V(\varphi)$ 可看成 $\Omega^1_{X/Y}\otimes\OO_{V(\varphi)}$ 的截面，它就是 $d_{X/Y}\varphi$ 在 $V(\varphi)$ 上的限制。因而可以考虑该截面的零点预概形，记为 $V(\varphi)_{\mathrm{sing}}$；由命题5，其底集正是 $\varphi$ 相对于 $Y$ 的奇异零点集。

（注意：若 $\psi$ 是预概形 $X$ 上有限型局部自由模 $E$ 的截面，则可以自然地定义 $\psi$ 的零点子预概形，例如由转置映射 $E^\vee\to\OO_X$ 的像理想定义。若 $E=\OO_X^n$、$\psi=(\psi_1,\ldots,\psi_n)$，则该理想为 $\sum_i\psi_i\OO_X$，定义的子预概形为 $V(\psi_1,\ldots,\psi_n)$。）

现在限制 $d^2\varphi$ 到 $V(\varphi)_{\mathrm{sing}}$，并注意
\[
 \operatorname{gr}^2(\PP^2_{X/Y})\simeq\Sym^2(\Omega^1_{X/Y}),
\]
便得到 $\Sym^2(\Omega^1_{X/Y})\otimes\OO_{V_{\mathrm{sing}}}$ 的典范截面 $M(\varphi)$。用域值点考察可知，这个截面正是定义1中决定二次型的截面（当 $X_k$ 由 $X/Y$ 经 $\Spec k\to Y$ 得到时）。\cmt{D17-S02}{底本此处两次使用$S$，本段仍以$Y$为相对底，故统一为$Y$；历史注号T08已经指出这一问题。}
由此可用该截面描述 $V(\varphi)_{\mathrm{sup\,sing}}$：将 $M(\varphi)$ 解释为同态
\[
 M(\varphi)':G_{X/Y}\otimes_{\OO_X}\OO_{V_{\mathrm{sing}}}
 \longrightarrow\Omega^1_{X/Y}\otimes_{\OO_X}\OO_{V_{\mathrm{sing}}},
\]
取这个同态不是同构的点集。特别地，这说明 $V(\varphi)_{\mathrm{sup\,sing}}$ 为闭集。假定 $X$ 在 $Y$ 上的相对维数恒为 $d$，可以引入
\[
 D(\varphi)=\det M(\varphi)
 \in\Gamma\left(V_{\mathrm{sing}},
       (\Omega^d_{X/Y})^{\otimes2}\otimes\OO_{V_{\mathrm{sing}}}\right)
\]
使描述更精确。于是也可用 $V(\varphi)_{\mathrm{sup\,sing}}$ 表示 $V(\varphi)_{\mathrm{sing}}$（从而也为 $X$）的闭子概形，由这个可逆模截面的消失定义；它有正确的底集。

\noindent\textbf{原稿编写提示：}最好把上述构造归纳为一个“命题6”。\cmt{P01}{所上传文件没有接着给出一个正式编号为6的命题陈述；本译文不补造。}

一般情形下，不能对 $V(\varphi)_{\mathrm{sing}}$ 与 $V(\varphi)_{\mathrm{sup\,sing}}$ 说得更精确。下面考察一个在某些应用中有意义的特殊情形。设 $Y$ 也对预概形 $S$ 光滑，为明确起见，设其相对维数恒为 $m$。为简便起见，将 $V(\varphi)_{\mathrm{sing}}$ 记为 $V'$；它由秩为 $d+1$ 的局部自由模 $\PP^1_{X/Y}$ 的截面 $d^1\varphi$ 的消失定义。再假设 $V'$ 对 $S$ 光滑，相对维数为\cmt{D17-S03}{底本写$(m+1)-(d+1)=m-1$，左边实际为$m-d$。由$X/Y$相对维数$d$及$Y/S$相对维数$m$，环境$X/S$的相对维数为$m+d$；原历史注T09已给出这一订正，本版将它移入正文。}
\[
 (m+d)-(d+1)=m-1.
\]


（注意，$V(\varphi)_{\mathrm{sing}}$ 与 $V(\varphi)_{\mathrm{sup\,sing}}$ 的记号有歧义，因为没有显示相对的预概形；此处约定相对底为 $Y$。还应注意，由假设可知，$\varphi$ 的每个相对于 $Y$ 的奇异零点，相对于 $S$ 都非奇异。）

在此情形下，可以写出下面的局部自由模层图。为使列图在默认版心内可读，设\cmt{图中文字}{以上 $A,B,C,D,R,W$ 只是本译稿为重排所加的缩写；原稿在图内展开写出各模。本图把原稿的 $V^1$ 统一显示为前文已定义的 $V'$，没有增加任何交换关系。$G$ 的含义沿用英译稿后面的编者疑问。}
\[
 A=(\PP^1_{X/Y})^\vee|_{V'},\quad
 B=\Omega^1_{X/S}|_{V'},\quad C=\Omega^1_{X/Y}|_{V'},
\]
\[
 D=\Omega^1_{Y/S}|_{V'},\quad R=\Omega^1_{V'/S},\quad
 W=\omega_{X/Y}^{-2}|_{V'}.
\]

\[
\begin{tikzcd}[row sep=large,column sep=large]
 &&0\\
 &&C\arrow[u]\\
 0\arrow[r]&A\arrow[r]\arrow[ur,"\mu"]&B\arrow[u]\arrow[r]&R\arrow[r]&0\\
 &\OO_{V'}\arrow[u]\arrow[r]&D\arrow[u]\arrow[ur,"\nu"']\\
 W\arrow[ur]&&0\arrow[u]
\end{tikzcd}
\]
竖直列来自光滑态射 $X\to Y$ 和 $Y\to S$ 的传递性正合列，再与 $\OO_{V'}$ 张量；由于列中各模均局部自由，张量后仍正合。横行是以下正合列的特例：在 $X/S$ 上给定局部自由模 $F$ 的截面 $\psi$，令 $W$ 为其零点概形，则有
\[
 F^\vee\otimes_{\OO_X}\OO_W\longrightarrow
 \Omega^1_{X/S}\otimes_{\OO_X}\OO_W\longrightarrow\Omega^1_{W/S}\longrightarrow0.
\]
若$X/S$光滑，则第一个同态在点$w\in W$的余域纤维上单射，当且仅当$W/S$在$w$处光滑，且相对维数为$X/S$的相对维数减去$F$的秩（这在当前所加的光滑性及维数假设下处处成立）。\cmt{D17-S04}{原文写的是模层同态“单射”，并把光滑的相对底印成$X$。模层单射本身不是Jacobian判据；正确条件为余域纤维满秩，相对底为$S$。首项也应限制到$W$。旧报告E06的尖点反例仍保留；本版将已有核查意见落实为正文的纤维秩判据。这里使用$F$的分量给出的局部方程，不能丢弃所列正确余维条件。}
这个正合列直接由第16段中的下列正合列得到（可以把上述版本列为推论）：\cmt{D17-S05}{底本末项重复印为$\Omega^1_{X/S}$；余法正合列的末项应为$\Omega^1_{W/S}$，本版落实旧注T10所列订正。}
\[
 J/J^2\longrightarrow\Omega^1_{X/S}\otimes\OO_W
      \longrightarrow\Omega^1_{W/S}\longrightarrow0.
\]


上述余域纤维上的单射条件，以及它所保证的列左端正合性，由Jacobian判据刻画。\cmt{D17-S06}{底本仅说“能够在最左边添上一个零”；本版承接前述纠正，将其限定为Jacobian判据的纤维秩条件，不再把抽象层单射等同于光滑性。}

注意，有典范同构
\[
 \PP^1_{X/Y}=\Omega^1_{X/Y}+\OO_X,
 \qquad (\PP^1_{X/Y})^\vee=G_{X/Y}+\OO_X
\]
（原稿中的$G$是拉长的字形）。\ednote{拉长的$G$是否就是切层？}
\cmt{D17-S07}{底本这里把三处相对下标写为$X/S$，但图中同态$\mu$的定义域为$(\PP^1_{X/Y})^\vee|_{V'}$，二次初始项来自$X/Y$；本版统一这些项为$X/Y$。$+$沿用原稿的直和记号，$G$沿用其切层记号，编者疑问照存。}
另一方面，可验证图中复合同态$\mu$在因子$\OO_{V'}$上为零，在因子$G_{X/Y}\otimes_{\OO_X}\OO_{V'}$上化为由$\Sym^2(\Omega^1_{X/Y})\otimes\OO_{V'}$的截面$M(\varphi)$得到的同态$M(\varphi)'$〔此处字迹不清〕。因此，在点$x\in V'$处，$M(\varphi)'$满射，当且仅当$\mu$在$x$处满射；由图中两条正合列，这又等价于$\nu$在$x$处满射，因为两者均表示$\Omega^1_{X/S}|_{V'}$的上述两个子模之和到该模的典范同态在$x$处满射。\cmt{D17-S08}{底本先把同态$\mu$印为$M$，再把$\nu$印为$V$；本版据原图订正（旧注T08）。这个满射性的线性代数等价不解决特征2下“退化二次型”的约定问题；不能由此声称本稿所有特征下的超奇异判据已独立验证。}
由此得到：

\egas{7}{命题}
在以上条件下（需重新列出），$V(\varphi)_{\mathrm{sup\,sing}}$ 的底集恰为 $V(\varphi)_{\mathrm{sing}}$ 中使态射
\[
 V(\varphi)_{\mathrm{sing}}\longrightarrow Y
\]
分歧的点集；这里源与靶都是 $S$ 上光滑的预概形，相对维数分别为 $m-1$ 与 $m$。

用前辈的说法（可把这一点写成备注），点 $x\in V(\varphi)$ 相对于 $Y$ 超奇异，当且仅当“它由至少两个重合的、无穷邻近的奇异点组成……”

可以而且必须把命题7精确化为子预概形之间的相等，而不仅是子集之间的相等。事实上，$V(\varphi)_{\mathrm{sup\,sing}}$ 已定义为 $X$ 的一个闭子预概形；同样也可自然地定义 $V'$ 的一个闭子概形，其底集为相对于 $Y$ 的分歧点集。只须表示局部自由模同态
\[
 Q=\Omega^1_{Y/S}|_{V'}\longrightarrow M=\Omega^1_{V'/S}
\]
不满射的点集。若两者的秩分别为$q,r$，这也是最高目标外幂同态
\[
 \bigwedge^r Q\longrightarrow\bigwedge^r M
\]
不满射的点集，亦即相应的典范截面
\[
 \mathscr{H}\!om(\bigwedge^r Q,\bigwedge^r M)
 \simeq(\bigwedge^r Q)^\vee\otimes\det M
\]
的零点集。\cmt{D17-S09}{底本破损式写为$\Lambda^1Q\to\Lambda\sqrt M$，并写$\Hom(\Lambda^1Q,\Lambda\sqrt M)\simeq(\approx)(\Lambda\sqrt Q)\otimes(\Lambda^rR)\otimes(\Lambda^1Q)^\vee$。这些字形在原PS渲染中确实存在，并非本次抽取造成。正确外幂为目标秩$r$的外幂，相应Hom模如正文；旧注T01已作此类型核对，本版将可确定的公式移入正文。}

原稿随即把它描述为这个显然截面的零点集，即一个闭零点子预概形的底集，并把它记为 $\operatorname{Ram}(V'/Y)$。我断言，这个子概形等于 $V(\varphi)_{\mathrm{sup\,sing}}$。这是关于上述图的一个简单练习：由 $\mu$ 的描述可知，$V(\varphi)_{\mathrm{sup\,sing}}$ 是对同态
\[
 P=(\PP^1_{X/Y})^\vee|_{V'}\longrightarrow
 S=\Omega^1_{X/Y}|_{V'}
\]
使用与刚才对 $Q\to R$ 所作相同的构造而定义的。

因此归结为以下一般情形。在赋环空间 $W$ 上给定秩 $m$ 的局部自由模 $M$，以及秩分别为 $p,q$ 的两个局部自由子模 $P,Q$，它们各自在$M$中局部为直和因子，\cmt{D17-S12}{这里要对$S=M/Q$及$R=M/P$使用局部自由模的行列式。原段开头只说局部自由子模；合并英文稿第66页紧接的论证明确使用$P,Q$局部为直和因子。本版把这一必需条件在使用商模之前列明。仅有子模自身局部自由不能保证商模局部自由。}使
\[
 p+q=m+1
\]
成立；对态射 $P\to M/Q=S$ 和 $Q\to M/P=R$ 使用上述构造，得到以下两个模的截面：\cmt{D17-S10}{底本把秩条件写为$p+1=m+1$，商模写为$M/W$。按余秩关系应为$p+q=m+1$，且$W$是赋环空间而非$M$的子模，故商模应为$M/Q$。本版落实旧注T01所列订正。}
\begin{align*}
 &P\otimes\det S\otimes(\det P)^{-1}\\
 &\hspace{1em}=P\otimes\det M\otimes(\det P)^{-1}\otimes(\det Q)^{-1},\\
 &Q\otimes\det M\otimes(\det P)^{-1}\otimes(\det Q)^{-1}.
\end{align*}

也可把这两个截面视为从
\[
 L=\det P\otimes\det Q\otimes(\det M)^{-1}
\]
分别到 $P,Q$ 的同态。\cmt{D17-S11}{底本最后一个因子写作$\det M$。由上面两个截面所属的模计算，应为$(\det M)^{-1}$；合并英文稿第66页紧接的图也明写此逆因子。本版据类型核对及该图落实旧注T01的订正。}
（注意：对局部自由模 $F$，用 $\det F$ 表示它的最高外幂，并使用如下事实：若
\[
 0\longrightarrow F'\longrightarrow F\longrightarrow F''\longrightarrow0
\]
是这种模的短正合列，则有典范同构 $\det F=\det F'\otimes\det F''$。）在这些条件下，有下面这个图的交换性——

\begingroup
\providecommand{\Oa}{\mathcal O}
\par\smallskip\noindent\textbf{载体接合。} 至此为六页短篇\texttt{EGAV1-2.ps}的末句；下接Piotr Blass与Joseph Blass翻译、整理的合并英文稿\texttt{EGAVaiello.pdf}第66--73页。译校注明法文扫描对读情况；不重复前段。\par
\[
\begin{tikzcd}[row sep=large,column sep=large]
P \arrow[r] & M \\
\det P\otimes\det Q\otimes(\det M)^{-1}=L
\arrow[u,"a"] \arrow[r,"b"] & Q \arrow[u]
\end{tikzcd}
\]
（可能相差一个符号；这或许取决于定义所用某些典范同构时采用的约定，而这些约定决定符号的选取……。）\transnote{图中的两条未标号箭头为子模到$M$的包含，$a:L\to P$、$b:L\to Q$的方向与底本相同。原文明确保留“possibly”“perhaps”及省略号，因此此处不擅自固定正负号，也不声称在未指定行列式同构约定时两条复合严格相等。即使相差$-1$，以下零点理想结论不变。底本图内明确写$\det M^{-1}$，即$(\det M)^{-1}$；这也与第十七版前段所作类型订正一致。}

因此，$V(a)$与$V(b)$是相同的零点子概形，因为$P$与$Q$在$M$中局部为直和因子。\transnote{底本原句为“the ideals $V(a)$ and $V(b)$ are identical”，混用了理想与零点对象的记号。这里将$V(a),V(b)$按上下文写为零点子概形；对应的零点理想相等，而不仅是底层点集相等。局部平凡化$L$后，包含$P\to M$有局部左逆，故一个$P$中向量在$P$内的坐标生成理想，与其在$M$中生成的坐标理想相同；$Q$同理。两条复合相差至多一个符号，遂给出相同理想。本说明是对原文理由的展开，并非另加新的全局符号约定。}
注意，在不属于$V(a)=V(b)$的点处，$L$恰可认同于$P$与$Q$在$M$中的交。\transnote{这里通过图中到$M$的复合来认同$L$的像；相差符号不会改变这个像。此断言只在共同零点子概形的补集上使用，不声称$L$在整个底空间上都等于交模。}

剩下的只是把上述图D所描述的、涉及$X/Y/S$和$\varphi$的特殊情形完全澄清，以便明确写出$a$和$b$。首先得到
\begin{align*}
\det P&=\omega_{X/Y}^{-1}\otimes\mathcal O_{V'},\\
\det Q&=\omega_{Y/S}\otimes\mathcal O_{V'},\\
\det M&=\omega_{X/S}\otimes\mathcal O_{V'}
       =\omega_{X/Y}\otimes\omega_{Y/S}\otimes\mathcal O_{V'}.
\end{align*}
（最后一个同构来自光滑态射$X\to Y$与$Y\to S$的$\Omega^1$传递性正合列。）
因此，\transnote{沿用前文$L=\det P\otimes\det Q\otimes(\det M)^{-1}$；代入以上三式，$\omega_{Y/S}$因子抵消，余下$\omega_{X/Y}^{-2}$。式中的等号按原文表示典范同构，而不是集合层面的相等。}
\[
L=\omega_{X/Y}^{-2}\otimes\mathcal O_{V'}.
\]
（注意：$\omega_{X/Y}$表示$\det\Omega^1_{X/Y}$。）这说明，同态$a$经过$(\mathcal P^1_{X/Y})^\vee\otimes\mathcal O_{V'}$的因子$\mathcal O_{V'}$分解；作为到这个因子的同态，它恰由前面已经提到的截面$D=\det M(\varphi)$确定（仍然相差一个有待确定的符号）。\transnote{底本此处印作$P^v_{X/Y}\otimes\mathcal O_{V'}$，省略了一阶上标，$v$表示对偶。根据前文$P=(\mathcal P^1_{X/Y})^\vee|_{V'}$及其直和因子，此处明确补出$1$并用$\vee$表示对偶；没有把它误作射影丛。底本的法语残片“en tant [Fr]”译为“作为”，不影响同态方向：这里是从$L$到该$\mathcal O_{V'}$因子。由于$L^{-1}=\omega_{X/Y}^{2}|_{V'}$，$D$的类型与这个同态相符。$M(\varphi)$是前文二次初始项的记号，不是本段环境模$M$本身。}
当然，利用引理中所宣布的交换性，就能得到同态$b$。\transnote{原文为“the commutativity announced in the lemma”，仍指上一批译出的图及其可能相差符号的约定；这里没有额外固定$a,b$各自的符号，也不自行补写下一段的正合列。}

因此，非超奇异点恰好是使
\[
W_{X/Y}^{-2}\otimes\Oa_{V'}\longrightarrow\Oa_{V'}
\]
成为同构的那些点。在这些点处，因子$\Oa_{V'}$便恰好可以与
$\mathcal P^1_{X/Y}\otimes\Oa_{V'}$及
$\Omega^1_{Y/S}\otimes\Oa_{V'}$在
$\Omega^1_{X/S}\otimes\Oa_{V'}$中的交相认同。
于是，在$V'$中由非超奇异点构成的开子集上，有正合列
\[
0\longrightarrow\Oa_{V'}\longrightarrow
\Omega^1_{Y/S}\times\Oa_{V'}\longrightarrow
\Omega^1_{V'/s}\longrightarrow0;
\]
因而还有正合列
\[
0\longrightarrow\Oa_{V'}\longrightarrow
\Omega^1_{X/S}\otimes\Oa_{V'}\longrightarrow
(\Omega^1_{X/Y}+\Omega^1_{Y/S})\otimes\Oa_{V'}
\longrightarrow0.
\]

应当用一个总结性的命题来概括整个情形，并以稍为一般的形式来表述，即取$X$上一个可逆模$K$的截面$\varphi$；这就要求把前面讨论中引入的若干模用$K$作扭曲。这确实就是超平面截面那一节所遇到的情形。

还要指出，若把$\varphi$换成$u\varphi$，其中$u$是一个单位，则
\[
V(\varphi)_{\mathrm{sing}}
\quad\text{和}\quad
V(\varphi)_{\mathrm{sup\,sing}}
\]
都不改变。事实上，$\mathrm d\varphi/V(\varphi)$及$M(\varphi)$都乘以$u$，而$D(\varphi)$乘以$u^M$，其中$M$是$X$相对于$Y$的维数。

\noindent\textbf{备注。}
在关于$X,Y,S,\varphi$的上述条件下，由于
\[
U=V(\varphi)_{\mathrm{sing}}-V(\varphi)_{\mathrm{sup\,sing}}
\]
在$Y$上非分歧，它在$Y$上的纤维是离散的。这意味着，对每个$y\in Y$，$X_y$中$\varphi_y$的每一个普通奇异零点在奇异零点集合中都是孤立的，也就是说，它在$V(\varphi_y)/k(y)$的非光滑点集合中是孤立的。

于是立即得出：若$d\ge 2$，它就是$V_y$在$k(y)$上的一个几何正规点（因为无论如何，$V_y$都是 Cohen--Macaulay 的）。\transnote{原文这一句附有“p.c.x.”，法文扫描为手写插入；此处保留缩写，不擅自展开。}

在定义1的条件下，可以更一般地证明：若$x$是一个普通奇异零点，则$x$是$V(\varphi)$的一个孤立奇点，其含义是，$x$在$V(\varphi)$中的每一个严格泛化点$x'$都是$V(\varphi)$的正则点。因此，若$V(\varphi)$的奇异点集是闭的，例如$X$是优良概形时，$x$也是该奇异点集的一个孤立点。（若求周全，应当把这个结论收入第四章，但第16节似乎并不是合适的位置。放在哪里呢？）

由此可见，若$X$在域$k$上光滑，而$x$是一个几何普通奇异零点，则$x$是$V(\varphi)$的非光滑点集合中的孤立点；因此，若$\dim_x(X)\ge 2$，它便是几何正规的。（但若$\dim_x(X)=1$，则必定不正规；例如可看典型情形$\varphi=xy$，其中$x,y$是仿射平面的坐标函数。）

下面给出一个与（D）类似的图，不过这里取可逆模$K$的一个截面$\varepsilon$。
\[
\begin{tikzcd}[column sep=0.8em,row sep=1.8em]
 & & 0 & & \\
 & & \Omega^1_{X/Y}\otimes\Oa_{V'} \arrow[u] & & \\
0 \arrow[r] &
 (\mathcal P^1_{X/Y}(K))^\vee\otimes\Oa_{V'}
 \arrow[r] \arrow[ur,"\mu"] &
 \Omega^1_{X/S}\otimes\Oa_{V'} \arrow[u] \arrow[r] &
 \Omega^1_{V'/S} \arrow[r] & 0 \\
 & K^{-1}\otimes\Oa_{V'} \arrow[u] \arrow[r] &
 \Omega^1_{Y/S}\otimes\Oa_{V'} \arrow[u] \arrow[ur,"\nu"'] & & \\
 & \omega_{X/Y}^{-2}\otimes K^{-1}\otimes\Oa_{V'} \arrow[u] &
 0 \arrow[u] & &
\end{tikzcd}
\]

\noindent\textbf{英译编者说明：} [以下部分被格罗滕迪克划去，英译者不确定原因。为求完整，仍将其转载；到划去部分结束之处，将标明“不再是划去内容”。]

\begingroup
\def\Sym{\operatorname{Sym}}
\def\gr{\operatorname{gr}}
\def\im{\operatorname{im}}
\def\Spec{\operatorname{Spec}}
\def\m{\mathfrak m}
\def\n{\mathfrak n}
\def\Oa{\mathcal O}
\def\Pp{\mathcal P}
\par\medskip\noindent\textbf{作者划去稿（以下至英译编者结束标记）。}\par
写到这个在一般性上堪称可观的地步，我却发现，前面这些展开仍过于特殊。与其考虑由$\mathcal O_X$的一个截面这一单个方程所定义的子预概形，不如考虑由$m$个截面所定义的子预概形。等价地，可以用一个同态
\[
E\longrightarrow\mathcal O_X
\]
定义它，其中$E$是秩为$m$的局部自由模。

是否从一开始就处理一般情形，留由你来决定，所依循的是从一般走向特殊的原则。\footnote{英译以方括号补出受话人迪厄多内；法文扫描此处只写“你”。}

先为带有增广$A\to A/I$的环$A$引入一套术语，以便也在层的情形中使用。若满足以下条件，就称此增广为\emph{正规的}（还是按照第19节采用的定义，称为拟正规的？）：
\begin{enumerate}
\item[(a)] $I/I^2$是有限型投射$A/I$-模；
\item[(b)] 同态
\[
\theta:\operatorname{Sym}(I/I^2)\longrightarrow\operatorname{gr}_I(A)
\]
是同构。
\end{enumerate}

若上面的（a）及下述（b$'$）成立，就称这一增广为\emph{普通二次型}的：
\begin{enumerate}
\item[(b$'$)] 前述典范同态有一个核，由$\theta^2$生成；这个核$K$是$\operatorname{Sym}^2(I/I^2)$的一个“非退化”的可逆子模。\footnote{此处照录有疑义的构造措辞；后一命题明写$K=\ker\theta^2$，详见译校说明。}
\end{enumerate}
所谓“非退化”，是指相应的同态
\[
K\otimes (I/I^2)^\vee\longrightarrow I/I^2
\]
为双射。这个同态由典范同态
\[
\operatorname{Sym}^2(I/I^2)\longrightarrow
\operatorname{Hom}((I/I^2)^\vee,I/I^2)
\]
导出。说它是双射，也等价于说它是满射，因为这里讨论的是$A/I$上在各点秩相同的投射模。换言之，局部取$K$的一个基，则$(I/I^2)^\vee$上相应的二次型非退化。其实，对$A/I$的各个素理想处的约化纤维检验后面这些条件就够了，乃至只在极大理想处检验也足够。

\noindent\textbf{命题8。}
增广$A\to A/I$为普通二次型的充要条件，是满足以下各条：
\begin{enumerate}
\item[$1^\circ$] 各$\gr_I^n(A)$都是有限型投射$A/I$-模。
\item[$2^\circ$] 若以$d$表示$\gr_I^1(A)=I/I^2$在$p\in\operatorname{Spec}(A/I)$处的秩，则$\gr_I^n(A)$在同一点处的秩等于（或者只需大于或等于）
\[
\binom{n+d}{d}-\binom{n-2+d}{d}.
\]
\item[$3^\circ$] （$1^\circ$及$2^\circ$已蕴含$K=\ker\theta^2$是可逆模，在$\Sym^2(I/I^2)$中局部为直和因子。）$\Sym^2(I/I^2)$的子模$K=\ker\theta^2$是非退化的。
\end{enumerate}

借助下述引理，证明是直接的。

\noindent\textbf{引理8.1。}
设$B$是一个环，$M$是一个有限型投射$B$-模，$\phi\in\Sym^2(M)$是一个“非退化”元素。那么，乘以$\phi$的同态
\[
\Sym^{[\text{次数待核}]}(M)\longrightarrow
\Sym^{[\text{次数待核}]}(M)
\]
是到某个直和因子子模的同构。因此，若$M$的秩为$d$，其余核是秩为
\[
\binom{n+d}{d}-\binom{n-2+d}{d}
\]
的投射模。\footnote{英译者已注明两处上标难辨，猜测可能涉及$n$和$n+2$。余核首项的上标英译印为$n+3$，法文第105页清楚为$n+d$，故本处据法文；此项仍保留原来的下标$d$。}

我们立刻归结到$B$是一个域的谱的情形，这时结论是平凡的。\footnote{法文也字面写$B$是一个域的谱，但$B$此前是环。语境中的归约是逐点取剩余域（其谱只有一点），并非将$B$等同于其谱。}（实际上，“非退化”只用来保证各个$\phi(x)$，$x\in\operatorname{Spec}(B)$，都不为零。用任意整数$k\ge1$代替$2$，而不提非退化，或许能更合理地表述本引理。）

\noindent\textbf{命题9。}
设环$A$带有正规增广$A\to A/I$，并且$I/I^2$是秩为$n$的投射$A/I$-模。设$E$是秩为$m$的投射$A$-模，$u:E\to I$是一个同态，因此有同态
\[
\bar u=u\otimes\operatorname{id}_{A/I}:
E\otimes_A A/I\longrightarrow I/I^2.
\]
令$B=A/u(E)A$，赋以增广$B\to B/J=A/I$，其中$J=I/u(E)A$。下列条件等价：
\begin{enumerate}
\item[(i)] 环$B$的增广$B\to B/J$为普通二次型，且秩为$n-(m-1)=n-m+1$。（增广的秩，按定义是投射模$J/J^2$的秩。）
\item[(ii)] $K=\ker\bar u$是可逆模，并且自然同态
\[
K\longrightarrow\Sym^2(J/J^2)
        =\Sym^2((I/I^2)/\im\bar u)
\]
是“非退化”的。
\end{enumerate}

如果记$V=I/I^2$、$W=J/J^2=V/\im\bar u$，也许更清楚。为得到（ii）中称为自然的那个同态，只须注意$K$无论如何都映入
\begin{align*}
I^2/(I^3+u(IE))
 &= (I^2/I^3)/\im(IE)\\
 &\simeq\Sym^2(I/I^2)/
       \im\bigl((I/I^2)\otimes_{A/I}(E\otimes_A A/I)\bigr),
\end{align*}
其中同态
\[
(I/I^2)\otimes_{A/I}(E\otimes_A A/I)
        \longrightarrow\Sym^2(I/I^2)
\]
是借助典范乘法同态
\[
(I/I^2)\otimes_{A/I}(I/I^2)\longrightarrow\Sym^2(I/I^2)
\]
定义的。作出这些说明之后，命题9的证明几乎是平凡的。

\noindent\textbf{推论10。}
设$A$是维数为$n$的正则局部环，$\m$为其极大理想，$k=A/\m$。设$E$为秩$m$的自由$A$-模，$u:E\to\m$为一同态。令
\[
\begin{gathered}B=A/u(E)A,\qquad \n=\m/u(E)A,\\ V=\m/\m^2,
\qquad W=\n/\n^2=V/\im\bar u,\end{gathered}
\]
其中$\n$是$B$的极大理想，$\bar u:E\otimes_A k\to V$由$u$诱导。下列条件等价：
\begin{enumerate}
\item[(i)] $B$的维数是$n-m$（也就是说，用$E$的一组基$(\phi_i)$来表述，这些$\phi_i$在$A$中的像组成一条$A$-正则序列），并且$\Spec(B)$的闭点是一个普通二次奇点。这里所谓普通二次奇点，是指增广$B\to B/\n$为普通二次型。
\item[(ii)] 存在$E$的一组基$\phi_1,\ldots,\phi_m$，使得：
\begin{enumerate}
\item[(a)] $A_{m-1}=A/(\sum_{i=1}^{m-1}A\,u(\phi_i))$是维数为$n-(m-1)$的正则环；
\item[(b)] $u(\phi_m)$在$A_{m-1}$中的像以$\Spec(A_{m-1})$的闭点为普通奇异零点（见定义1）。
\end{enumerate}
\item[(ii bis)] $\im\bar u$的秩为$m-1$；对$E$的任意一组基$(\phi_i)_{1\le i\le m}$，只要$\phi_i$（$1\le i\le m-1$）在$V$中的像线性无关，（ii）的（a）和（b）就成立。
\item[(iii)] $\bar u:E\otimes_A k\to V$的核$K$是维数为一的向量空间；对任意$\phi\in K\setminus\{0\}$，由$\phi$所确定的$\Sym^2(W)$中的元素是一个非退化二次型。
\end{enumerate}

\noindent〔英译者按：划去的部分到此结束。以下不再是划去内容。〕

\par\medskip\noindent\textbf{零点的术语与奇异子预概形（恢复未划去正文）。}\par
请注意，（ii）和（ii bis）的类似条件，本可以提前放到命题9中。应当引入如下术语。设$X$为正则预概形，$E$是$X$上的局部自由模，$u:E\to\Oa_X$是同态，即$E^\vee$的一个截面。对
$x\in V(u)=V(u(E)\Oa_X)$，若$u$是$\Oa_X$-正则的，即
\[
\dim\Oa_{V(u),x}=\dim\Oa_{X,x}-\operatorname{rank}_x E,
\]
并且$x$是$V(u)$的正则点，就称$x$为$u$的\emph{非奇异零点}。取$x$的一个邻域上$E$的一组基$\phi_1,\ldots,\phi_m$，这也可以表述为：各$u(\phi_i)$是$\Oa_{X,x}$的某个正则参数系的一部分。否则，称$x$为\emph{奇异零点}。

若$x$满足推论10中的等价条件，就称$x$为\emph{普通奇异零点}（因此它确实是奇异零点）；这里自然取$A=\Oa_{X,x}$，其余记号相应解释。若$x$是奇异零点，却不是普通奇异零点，就称为\emph{超奇异零点}。因此，非超奇异零点就是非奇异零点或普通奇异零点。

最后，当$X$在域$k$上光滑时，可以用显然的方式引入这些概念的几何版本。注意，撇开维数条件，也就是$\Oa_X$-正则性条件不谈，$x$的性质实际上可从它作为$V(u)$上的点的性质看出，即可从局部环$\Oa_{V(u),x}$看出（必要时将它视为$k$-代数）。在基域为$k$的情形，若$x$是$k$-有理点，则相对于$k$的概念与绝对概念一致，这可像$E$为可逆模时那样看出。最后，若$X$在预概形$S$上光滑，而$u:E\to\Oa_X$仍为同态，$E$为有限型局部自由模，就逐纤维地定义相对于$S$的对应概念。注意，若$x$相对于$S$非超奇异，则$V(u)$在$x$处在$S$上平坦（因为$u$相对于$S$是$\Oa_X$-正则的……）。\footnote{平坦据法文第108页清楚的“plat”；英译此处作“is that”。英译者另注基底字母曾改写，可能是$Y$；本段按逐纤维讨论所用的同一基底记为$S$。}

当然（在定义正式写定以后，至少应把这点作为备注指出），也完全可以从局部自由模$E$的截面$u$出发，把$u$理解为同态$E^\vee\to\Oa_X$，即可应用前述定义。

当$X$在$Y$上光滑，给定同态$u:E\to\Oa_X$且$E$为秩$m$的局部自由模时，$u$相对于$Y$的奇异零点集是$V(u)=V^0$中一个确定的闭子预概形$V(u)^{\mathrm{sing}}=V^1$的底层集合。构造如下：取$\mathrm d^1u$，即$\mathcal P^1_{X/Y}(E^\vee)$的一个截面。它在$V^0$上的限制可以解释为同态
\[
E\otimes\Oa_{V^0}\longrightarrow
\Omega^1_{X/Y}\otimes\Oa_{V^0}.
\]
于是$V(u)^{\mathrm{sing}}=V^1$表示相应同态
\[
(\bigwedge^m E)\otimes\Oa_{V^0}\longrightarrow
\Omega^m_{X/Y}\otimes\Oa_{V^0}
\]
的零点预概形，也就是以下线丛扭曲的微分模中一个确定截面的零点预概形：
\[
\Omega^m_{X/Y}\otimes(\det E)^{-1}\otimes\Oa_{V^0}.
\]

要用微分方法刻画相对于$Y$的普通奇异零点，只需对从$u$导出的同态
\[
\Pp^\infty_{X/Y}(E)|_{V^1}\longrightarrow
\Pp^\infty_{X/Y}|_{V^1}
\]
表达命题9的条件。实际上，这又化为一个“二阶”条件，只涉及相应的同态
\[
\Pp^2_{X/Y}(E)|_{V^1}\longrightarrow
\Pp^2_{X/Y}|_{V^1}.
\]

首先应表达同态
\[
(*)\qquad E|_{V^1}\longrightarrow\Omega^1_{X/Y}|_{V^1}
\]
的核$K$在$x$处可逆且为直和因子。换言之，该同态在$x$处的秩（因$x\in V^1$而至多为$m-1$）恰好为$m-1$。\footnote{“至多”是由上一批顶次外幂的零概形定义补明的语义；英、法在括号内都直接写$m-1$。}（注意：局部自由模间同态在一点的秩，按定义就是它在$k(x)$上的相应向量空间同态的秩。）为表达这一点，取同态
\[
(\bigwedge^{m-1}E)|_{V^1}\longrightarrow
\Omega^{m-1}_{X/Y}|_{V^1}
\]
的零概形$V''\subset V^1$，并只须写$x\notin V''$。\footnote{法文第109页明确为$x\notin V''$；英译排成$x\in V''$。以下用法文可确定的$V^1$，不将英译中的撇号读成另一层概形。}

$V''$的定义确保$K|_{V^1\setminus V''}$是可逆模，且$(*)$的余核$W$在$V^1\setminus V''$上局部自由，秩为$d-(m-1)$（其中$d$是$X$在$Y$上的相对维数，等于$\Omega^1_{X/Y}$的秩）。\footnote{此处法文维数叙述有印字$S$，相邻相对微分的底符有手改。英译者也指出这里的$S$应为$Y$。本段指同一相对构造。}应用命题9中说明的构造，得到同态
\[
M:K\longrightarrow\Sym^2(W).
\]
只剩下表达它在$x$处是“非退化”的。为此，引入判别式
\[
D\in\Gamma\bigl(V^1\setminus V'',\,
 K^{-d+m-1}\otimes(\det W)^2\bigr)
\]
（指数表示可逆模的张量幂）。\footnote{法文及英译原式均为$K^{-n+m-1}$。因$W$的秩为$d-m+1$，取极化映射的行列式得正文的幂次；不据此推断作者原来如何区分$n$与$d$。}令$V^2=V(D)$。它是$V^1\setminus V''$的闭子预概形。于是相对于$Y$的普通奇异零点集就是
\[
V^1\setminus(V''\cup V^2),
\]
也就是说，相对超奇异零点的集合为$V''\cup V^2$。\footnote{末句源文的两个记号之间未清楚印出并号。此处用前一句的补集关系确定并集；英译者另有括号疑问。这里仅陈述底层点集。}

引入$V''$时，我在断言$W$与$K$在$V^1\setminus V''$上局部自由时用了一个一般事实；现在将它单列为引理（其中$V^1\setminus V''$变成$X$）。

\noindent\textbf{引理。}
设$u:E\to F$是预概形$X$上局部自由模之间的一个同态，$m>0$为整数，并假设$\bigwedge^m u=0$。那么以下条件等价：$\bigwedge^{m-1}u$在任一点都不消失（当然按各点的约化纤维理解）；$u(E)$是$F$的一个局部直和因子，并局部自由且秩为$m-1$；$u$在$X$的每一点秩都为$m-1$。

宜把这些构造概括为一个推广命题6的命题。

\enlargethispage{4\baselineskip}
还须看一看怎样同样推广命题7。为此，假设$Y$本身在预概形$S$上光滑，相对维数为$N$。于是$X$在$S$上光滑，相对维数为$N+d$，其中$d$是$X$在$Y$上的相对维数。\footnote{英译把$d$说成$Y$的相对维数；法文跨行读作“其中$d$是〔$X$〕在$Y$上的相对维数”，对象在扫描右缘不清，但由$N$已为$Y/S$维数及前文$d=\operatorname{rank}\Omega^1_{X/Y}$可确定此订正。}
沿用前面的记号，再假设$V''=\varnothing$，于是$W$是秩为$d-m+1$的局部自由模（当$m=1$时总是如此！）。

最后，假设$X$的闭子预概形$V^0$和$V^1$都在$S$上光滑，并且在每一点都具有给定$d,N,m$时所设定的最小相对维数，也就是分别具有相对维数
\[
(N+d)-m,\qquad (N+d)-m-(d-m+1)=N-1.
\]
（从行文上说，最好直接给出这些数，或许改用逐纤维余维表述：$V^0$在$X$中的余维为$m$，$V^1$在$V^0$中的余维为$d-m+1$；再在备注中用“泛型”最小维数原则说明这样选取的理由。）\footnote{法110手写补明是$V^0$和$V^1$，英译重复印成$V^0$。光滑性及预期余维等式是独立假设，不由下文给出的行列式理想高度上界推出。}

设定这些条件之后，应当验证：$u$相对于$Y$的超奇异点集$V^2$恰是$V^1\to Y$发生分歧的点集，并把这一点精确表述为$X$的闭子预概形的恒等式：
\[
V^2=\text{$V^1\to Y$的分歧子预概形}.
\]
我没有详细做过这个验证，但相信基本上可以沿用命题7之后的那种分解方法。

\medskip
\noindent\textbf{附注。}
设$u:E\to F$是正则预概形$X$上局部自由模之间的同态，$E,F$的秩分别为$m,d$（$m\le d$）。若$Z$由$u$的秩至多为$m-1$的点组成，则$Z$的每个不可约分支在$X$中的余维至多为$d-m+1$。\footnote{法111和英73都写余维$\ge d-m+1$。该方向一般错误：在$X=\operatorname{Spec}k[t]$上取$\OO_X\to\OO_X^2$，$1\mapsto(t,0)$，则$Z=V(t)$的余维是$1$，而$d-m+1=2$。正文改为行列式理想高度定理给出的上界。参见 Eisenbud--Evans, \emph{A generalized principal ideal theorem}, Cor.~2.3（Eagon--Macaulay--Northcott），\href{https://doi.org/10.1017/S0027763000024740}{DOI 10.1017/S0027763000024740}；在每个分支的泛点局部化，取子式次数$m$。英译把这则附注提前，法111置于分歧段之后，本稿从法文次序。}
\endgroup

\par\medskip\noindent\textbf{本次新增译校附记（以下并非作者原文）。}\par

\paragraph{符号与法文核对。}
上述显示式保留英译第67页的$W$、$\times$、小写$s$及加号；法文扫描第102页对应处分别可读为$\omega$、$\otimes$及大写$S$。结合上文，应理解为线丛$\omega_{X/Y}^{-2}|_{V'}$的同态及限制后的微分模。这里来自$Y$的模先拉回，再限制。

交模句的$\mathcal P^1_{X/Y}$在英译与法文扫描该句中均未见对偶标记；但前文设定为$P=(\mathcal P^1_{X/Y})^\vee|_{V'}$，且法文第102页首承接句明确带对偶记号。交模分析应使用前文的$P$，不能不加说明地将无对偶的主部模视为环境模的子模。

\paragraph{第二条正合列的秩问题。}
法文扫描第102页第二条正合列的目标同样写作
$(\Omega^1_{X/Y}+\Omega^1_{Y/S})\otimes\Oa_{V'}$，所以这里不是仅凭英译版式猜测的疑点。若把加号解释为未经商的两个微分模的外直和，设$X/Y$、$Y/S$的相对维数分别为$d,m$，则中项秩为$d+m$，左端核秩为$1$，目标却仍为$d+m$，不可能正合。

\paragraph{按前图得到的条件性解释。}
令$U$为非超奇异开集，取
\[
E=\Omega^1_{X/S}|_U,\quad
P=(\mathcal P^1_{X/Y})^\vee|_U,\quad
Q=\Omega^1_{Y/S}|_U.
\]
按照前图的嵌入与商映射，在$U$上共同线
$K_0=P\cap Q\simeq\Oa_U$，且$P+Q=E$。因此自然映射
\[
(P/K_0)\oplus(Q/K_0)\longrightarrow E/K_0,
\qquad(\bar p,\bar q)\longmapsto\overline{p+q}
\]
为同构：和条件给满射；若$p+q\in K_0$，则$p\in P\cap Q=K_0$，进而$q\in K_0$，故核为零。

又由前图的商模辨认，
\[
P/K_0\simeq E/Q\simeq\Omega^1_{X/Y}|_U,
\qquad
Q/K_0\simeq E/P\simeq\Omega^1_{V'/S}|_U.
\]
所以在这些条件下，类型和秩一致的正合列为
\[
0\longrightarrow\Oa_U\longrightarrow E\longrightarrow
\Omega^1_{X/Y}|_U\oplus\Omega^1_{V'/S}|_U\longrightarrow0.
\]
两个右侧因子的秩分别为$d$和$m-1$。必须先将$Q$商去共同线；不能直接把原有的$\Omega^1_{Y/S}$拿来相加。以上是基于前图和和、交条件的校勘推断，并非法文原文已经写出的订正式，故未用它替换译文。

\paragraph{原文对读与记号。}
法文扫描第102页的“Il faudrait résumer”是尚待总结的语气，不是已给出某个编号命题。该段关于单位变换的后半句是手写补入，英译已将其收入正文。法文“tordre par $K$”对应“用$K$作扭曲”；本段并未逐一列出应当张量的幂次，不能据此把后面的图擅自补齐。

英译的$\mathrm d\varphi/V(\varphi)$表示将微分限制到零概形，并非将微分除以一个概形。$M(\varphi)$沿用前文二次初始项的矩阵记号，而同句单独出现的大写$M$表示相对维数；两者不能混同。
\paragraph{单位变换的适用位置。}
下面只是说明原句的类型与限制条件，不替代原文。
局部把截面写成函数，并在相同局部标架下比较。由乘积法则
\[
\mathrm d(u\varphi)=u\,\mathrm d\varphi+\varphi\,\mathrm du,
\]
限制到$V(\varphi)=V(u\varphi)$后，第二项消失，故限制后的微分确实乘以$u$。在临界零概形上再令$\varphi$及其相对一阶微分为零，二次初始项中含这些因子的交叉项也消失，因而其矩阵乘以$u$，行列式乘以$u^M$。

这不是说在整个$X$上未经限制的二阶导数矩阵恒满足同一缩放等式。上述单位因子不改变相应的零理想，说明本段所说的奇异及超奇异零概形不变。这里仍沿用底本的定义和术语；特征$2$下二次型光滑性与双线性判别式的区别，须保留主稿既有校勘，不能以本段的单位不变性消除该区别。

\paragraph{措辞及范围。}
法文第102页在孤立性句中手写补入“即”的意思，与英译一致。第103页所记“每个泛化$x'$”按此处论旨指异于$x$的泛化；若包含$x$本身，就会同时断言该奇点正则。英译中疑为排字错误的“generalization of $x^1$ of $x$”据法文读为“$x$的泛化$x'$”。定义1段是原文提出的进一步结论，不是原文已给出证明。此处保留历史术语“普通奇异零点”；特征$2$下它与双线性判别式的关系仍应参照主稿已有校勘。

\paragraph{正规性的维数疑点。}
英译第67--68页与法文第102--103页均保留正文所写的阈值；尤其末段法文也写$\dim_x(X)\ge2$与$\dim_x(X)=1$。不在译文中替作者改动。这里须分清环境$X$和超曲面$V(\varphi)$的维数：若在几何纤维的局部环上，$V(\varphi)$为 Cohen--Macaulay，并且除闭点外局部均正则，则 Cohen--Macaulay 给出$S_2$；当该局部环维数至少为$2$时，所有高度至多$1$的点都不是闭点，故满足$R_1$，由 Serre 正规性判据得正规。这个论证的维数条件是加在超曲面局部环上，不是直接加在环境上。

原文自己的例子说明不能把末段的环境维数阈值照字面作为无条件结论：在仿射平面的原点，环境维数是$2$，但$xy=0$的局部环有两个分支，非整环，因而不是正规局部环；零概形在此处的维数是$1$。首处$d$究竟取零概形维数还是环境相对维数，须按前文具体定义核对，不因末段疑点而擅自代换。此处的维数辨析是译校推断，不是原文订正版。

\paragraph{图的对读。}
法文第103页引图的手写截面符号看起来作$\varphi$；本批重绘沿用英译的$\varepsilon$，不因此另设两个截面。图中$v$为对偶符号，重绘用$\vee$。$\mu$是先由左中项入环境微分模、再映至相对微分模的复合；$\nu$由下中项经环境微分模映至右端。图中的顶端和底端零对象属于环境模的竖直正合列，并非另为斜箭头断言单射或满射。

\paragraph{左下角的指数。}
重绘照英译保留$K^{-1}$；法文第103页对应手写指数有涂抹，不能据此确认完整字样。下面只作类型核对。若仍沿用前文
$L=\det P\otimes\det Q\otimes(\det E)^{-1}$，且$P$为秩$d+1$的相对一阶主部模$\mathcal P^1_{X/Y}(K)$之对偶，那么在主部模正合列
\[
0\to\Omega^1_{X/Y}\otimes K
\to\mathcal P^1_{X/Y}(K)\to K\to0
\]
的约定下，
$\det P\simeq \omega_{X/Y}^{-1}\otimes K^{-(d+1)}$。
取$Q=\Omega^1_{Y/S}|_{V'}$和$E=\Omega^1_{X/S}|_{V'}$，将这些线丛拉回并限制至$V'$，传递性行列式关系遂给出
\[
L\simeq
(\omega_{X/Y}^{-2}\otimes K^{-(d+1)})|_{V'}.
\]
因此，若左下角意在表示这条$L$，则指向$K^{-1}|_{V'}$的箭头相应属于
$\operatorname{Hom}(L,K^{-1}|_{V'})
\simeq(\omega_{X/Y}^{2}\otimes K^{d})|_{V'}$，
其$K$幂次与上一批单位变换的$u^d$相容。这不是对模糊手写指数的辨读，也不是在译文中静默更换英译图的标签。

\paragraph{划去标记。}
方括号是英译者说明，不是作者正文。法文第104页从下一段已可见贯页划线；第一辑止于这条说明；第二辑已查定整段划去稿的最终边界，并恢复其后的非划去正文。

\begingroup
\def\Sym{\operatorname{Sym}}
\def\gr{\operatorname{gr}}
\def\im{\operatorname{im}}
\def\Spec{\operatorname{Spec}}
\def\m{\mathfrak m}
\def\n{\mathfrak n}
\def\Oa{\mathcal O}
\def\Pp{\mathcal P}
\par\medskip\noindent\textbf{第二辑新增译校附记（以下并非作者原文）。}\par

\par\medskip\noindent\textbf{56：作者划去稿（续）。}\par
\paragraph{（a）不等于（a）和（b）。}
英译第69页的“the above”可能让人以为要求上一节的两个条件。法文第104页却明确写“a) ci-dessus”，另列手写编号（b$'$）；因此根据法文只纳入（a）。正规增广要求$\theta$是同构，普通二次型增广另要求二次关系，二者不是累加条件。英译把同态字形排作$O$，对应法文的$\theta$。

\paragraph{核及生成所在的次数。}
法文第104页“由$\theta^2$生成”的字面措辞与后半句的简写并不能直接作为类型正确的定义：$\theta^2$是映射，不是生成元；全次数的齐次理想也不是二次分量的可逆子模。法文第105页命题8第3项明确写$K=\ker\theta^2$。据此可推断意旨是二次分量$K=\ker(\theta^2:\operatorname{Sym}^2(I/I^2)\to\operatorname{gr}_I^2(A))$为可逆子模，并在整个对称代数中生成$\ker\theta$作为齐次理想。后者是据上下文提出的解释，不是说扫描清楚写出了这一完整现代定义。

\paragraph{局部条件和划线。}
上面的非退化按正文实际使用的同态判别；在特征$2$不能自行改成射影二次超曲面光滑。对等秩有限投射模，满射与同构的等价性可在各局部环上用行列式检验。英译者的删除段终止标记在英译第71页推论10（iii）之后。法文第104--107页可见粗划线；第107页的粗线穿过推论10（iii）并在下一段前结束，下一段另有边栏记号，不能与贯页划去线混同。本批未翻译命题8及其后的划去文字。

\par\medskip\noindent\textbf{57：作者划去稿（续）。}\par
\paragraph{命题8的秩障碍。}
法文明确补出$d$为一次分量的秩，且两项的二项式下标均为$d$，也保留“或只需$\ge$”。不能只将这些缺陷归于英译。令$n=2$，原式化为$\binom{d+2}{d}-1=d(d+3)/2$；而$\theta^2$为满射，故二次分量秩至多$\operatorname{rk}\Sym^2(I/I^2)=d(d+1)/2$。当$d\ge1$，前者比后者多$d$，连下界也不可能。由此不能沿用原式推出$3^\circ$括号中的结论。

\paragraph{引理的可复原内容。}
法文两次数上标有叠写，英译者也标明难辨。按目标处于$n$次来重建，乘二次式的映射应为
$\Sym^{n-2}(M)\to\Sym^n(M)$（$n<2$时源为零）；也可把源次数记为$n$，则目标为$n+2$。这不是辨读出法文最终字样。若$\phi(x)$在每个剩余域上非零，纤维的对称代数为多项式环，乘$\phi(x)$为单射。将有限投射模局部平凡化后，源秩阶的某个子式在每一点附近可逆，故原映射局部裂单，余核有限局部自由，从而是有限投射模。固定正秩$d$时，所得余核秩为
\[
h_n=\binom{n+d-1}{d-1}-\binom{n-2+d-1}{d-1},
\]
其中负次数的对称幂按零计。若改用次数$k$，同一论证只要求每个纤维中的$k$次式非零。

\paragraph{即使改正计数，下界也不能独自保证有二次关系。}
设$A=k[T_1,\dots,T_d]$，$I=(T_1,\dots,T_d)$，$d\ge1$。则$\gr_I(A)=\Sym_k(I/I^2)$；每个分量均有限维且维数至少为$h_n$，但$K=0$，并非可逆模。这反驳“$1^\circ$和下界已推出$K$可逆”的说法，不反驳另行假定$K$为可逆且非退化后的判别。后者在有限投射及各纤维中$K$的局部生成元非零的条件下，先由上述引理给出$\Sym(I/I^2)/(K)$各次秩为$h_n$，再与其到$\gr_I(A)$的各次满射比较；此时$\gr_I^n(A)$的秩下界$h_n$足以使满射成为同构。这些是数学重建，不是作者已经写出的替换命题。

\par\medskip\noindent\textbf{58：作者划去稿（续）。}\par
\paragraph{英译的段次与商式。}
英译第70页把自然同态解释的一部分插在（i）与（ii）之间，把二次商式的最后一段割裂，并在等号后漏排商号。法文第106页先完整列（i）、（ii），后给连续商式；译文按这个可见顺序还原。法文第一商的$u(IE)$清楚，末项靠右纸缘不全。整理后的末项商掉的是所显示乘法映射的像；这是由紧随原文的映射及对称代数的商性质确认的，不是按缺字猜补。$\im(IE)$是$u(IE)$在$I^2/I^3$中的像，不是$IE$本身。

\paragraph{自然二次类与提升的选择。}
令$C=A/I$，$H=\im\bar u$，$W=V/H$。取$k\in K\subset E/IE$，选取其提升$e\in E$。因$\bar u(k)=0$，有$u(e)\in I^2$。若换另一提升，则差属于$IE$，使$u(e)$的变化落入$u(IE)$；故$k\mapsto[u(e)]$在$I^2/(I^3+u(IE))$中定义良好，且是$C$-线性的。由正规增广的二次识别$I^2/I^3\simeq\Sym_C^2(V)$，$u(IE)$的像恰为$V\cdot H$，即由$v\,\bar u(\bar e)$生成的子模。于是该目标是
\[
\Sym_C^2(V)/(V\cdot H)\simeq\Sym_C^2(W).
\]
这一步不先要求$W$投射；它使用的是正规增广条件以及对称代数对商模的函子性质。

\paragraph{对等价性的有限重建。}
原文只说证明几乎平凡。在$W$是秩$n-m+1$的有限投射$C$-模、且（ii）的自然二次类在每个纤维非零的条件下，可给出如下局部论证。此时$H$在$V$中局部为直和因子，秩$m-1$；满射$E\otimes_A C\twoheadrightarrow H$亦局部裂开。沿$\operatorname{Spec}C$局部化并在$A$中提升相应基，$m-1$个$u$的初始项是$V$的独立线性变量，余下的生成元$f$在消去这些方程后属于平方滤层，二次初始项为上述自然二次类的局部生成元。

逐次消去每个一次变量时，以及最后消去$f$时，使用同一事实：设$g\in I^r$且其$r$次初始项为非零因子。若$ag\in I^N$、$N\ge r$，逐次察看$a$在次数$0,1,\dots,N-r-1$的初始类，得$a\in I^{N-r}$。因此商环的分次环就是原分次环除以该初始项。一次变量是非零因子；最后二次式由上一批引理的逐纤维非零论证，在各次乘法都是单射。所得分次环为$\Sym_C(W)$除以该二次关系。逆向若假设（i），$W$已有限投射且秩指定，故上述局部消去一次变量可做；剩下单个$f$属于平方滤层，二次核由它在平方滤层中的类生成。这个核按（i）是可逆的；自然映射由可逆模$K$满射到它，因而为同构。

这里不把“$K$本身可逆”自动当成$\bar u$在每个剩余域均秩$m-1$的证明。例如$k[t]^2\to k[t]^2$、$(a,b)\mapsto(ta,0)$的核自由秩一，但在$t=0$处纤维映射为零，余核有挠。此例只反驳该中间推理，不声称满足完整（ii）。以上重建将投射及纤维条件显式列出，不把独立假设冒充原稿已有的证明。

\par\medskip\noindent\textbf{59：作者划去稿（续至结束）。}\par
\paragraph{可由法文确定的订正。}
法文第106页明确写“anneau local régulier”，因此补回英译漏掉的“正则”。英译在同句也混淆了秩$m$、极大理想$\m,\n$、剩余域$k$和核$K$。译文按两页法文及前一命题的记号，将“像”统一作$\im\bar u$；（ii）（a）中商掉的是由相应元素生成的理想。法文第107页明确给出（ii bis）的像秩和（iii）的核维数。这里没有用主稿后来讨论的特征$2$术语替换历史原文。

\paragraph{只对英译漏项的反例。}
以$k=\mathbb Q$，取$A=k[[x,y]]/(x^3)$，$E=A$，$u(1)=x^2+y^2$。这是维数为$1$的非正则局部环：$y$给出一维余环方向，而$x,y$在$\m/\m^2$中独立，嵌入维数为$2$。由于关系$x^3$到三次才出现，$\Sym_k^2(V)\to\m^2/\m^3$在二次仍是同构，故二次类在此例可直接辨认。$\bar u=0$且$K=k$，非零元所给二次型均为$x^2+y^2$的非零倍数，在$\mathbb Q$上非退化，故（iii）成立。可是$m=1$时（ii）（a）的空组商环就是$A$，它不正则，故（ii）不成立。这只说明不能把英文漏掉“正则”的文字当作一般局部环的命题，不是法文命题的反例。

\paragraph{划去范围。}
法文第107页贯页粗线穿过（iii）并止于其末附近；下一段另有边栏记号。英译者的结束标记单列，不计入作者正文，也不表示该推论已收入后来定稿。

\par\medskip\noindent\textbf{60：零点的术语与奇异子预概形（续）。}\par
\paragraph{底本与基底。}
法文第107页结束标记后开头另有边栏记号，但已无此前跨页斜线。法文第108页基底处可见印字$S$与手改字样交替；英译者也标注犹疑。本稿在逐纤维定义一段统一用$S$，在下一段微分构造统一用$Y$，并不宣称每处手改字母的笔画皆可确定。地位相同的底基更名不改变公式；不能同时把不同基底的相对微分混为一个模。

\paragraph{外幂的类型。}
英译第71页把最后的源排作$E|_{V^0}$，法文第108页末明确在$E$前写$\bigwedge^m$。取前一同态的第$m$外幂，因此源为$\det E|_{V^0}$；$\Omega^m_{X/Y}$指$\bigwedge^m\Omega^1_{X/Y}$，本身不一定为线丛。局部选基后，该截面的分量为微分矩阵的$m$阶子式，故其零点预概形是所需的秩下降闭子预概形。

\paragraph{相对平坦的条件。}
非超奇异包括前文两类，均带纤维上的预期余维条件。在光滑环境的正则纤维局部环中，满足该条件的局部$m$个方程构成正则序列；光滑环境本身对基底平坦。在通常的局部平坦性判据适用时，逐个商去这些在纤维上正则的方程，得到零点预概形在该点相对平坦。此处是对原文简述的解释，不是从某个二次型在单一纤维上非退化、却无维数条件直接推出平坦。

\par\medskip\noindent\textbf{61：普通奇异零点的微分刻画（续）。}\par
\paragraph{概形层面恒秩。}
在$V^1$上，$m$阶子式全部在结构层中为零，而不只是在每个剩余域中为零。在$V^1\setminus V''$上，局部某个$(m-1)$阶子式是单位。用它作主块，初等可逆行列变换将矩阵化为$\operatorname{diag}(I_{m-1},C)$的块形；$C$的每个条目都是变换后一个$m$阶子式，故$C=0$。因此即便基底非约化，核、像、余核仍局部自由；源模的秩为$m$时，核秩为一。引理的两种纤维秩表述直接等价；相应的局部直和形式由这一步给出，反向由局部直和形式在纤维上取秩得到。仅有逐点秩而无$\bigwedge^m u$作为层同态为零，不能这样消去幂零项。

\paragraph{判别式的类型。}
令$r=\operatorname{rk}W=d-m+1$。由$M$的极化得到$K\otimes W^\vee\to W$。取行列式，其值线丛是
\[
\operatorname{Hom}(\det(K\otimes W^\vee),\det W)
 \simeq K^{-r}\otimes(\det W)^2.
\]
因此按本段的$d$记号，指数应核算为$-d+m-1$；法文与英译可读的$-n+m-1$保留在脚注。这个行列式检测极化映射是否退化；在特征$2$下，不能未经另证将它等同于二次型的射影零超曲面是否光滑，这与主稿已有校勘一致。法文第112页（英译第73页）后记专门回到此问题，举$x^2+yz$为例并提出术语调整；本批未译该后记。

\paragraph{分层区别。}
$V^2$在开层$V^1\setminus V''$中闭，不必单独在$V^1$中闭；$V''\cup V^2$则是普通层在$V^1$内的补集。它不包括$V^0\setminus V^1$的非奇异零点。原文只在此辨认集合，没有指定合并后的闭子概形结构。

\par\medskip\noindent\textbf{62：光滑塔中的分歧（续）。}\par
\paragraph{分歧子概形的局部核对。}
以下补证不是原稿的证明。取$V^1$上的一点。由于$V''=\varnothing$，在该点至少有一项$(m-1)$阶相对雅可比子式为单位。平凡化$E$并重排局部方程$u_1,\dots,u_m$后，可将$u_1,\dots,u_{m-1}$连同$r=d-m+1$个函数$t_1,\dots,t_r$取作$X/Y$的相对 étale 坐标：这些函数的相对微分在该点成基，对应到$\mathbb A^d_Y$的态射局部为 étale。令
\[
T=V(u_1,\dots,u_{m-1}),\qquad f=u_m|_T.
\]
则$T/Y$光滑，相对坐标为$t_i$，而$V^0=V(f)$是在$T$中所取的零概形。因前$m-1$行可化为单位块，顶次外幂的子式理想在$T$上恰给出余下相对偏导，故$V^1$在$T$内由理想
\[
(f,\partial_1 f,\dots,\partial_r f)
\]
定义。这是理想的等式，不仅是点集等式。

对这个商取相对微分。$\mathrm df$在$V^1$上为零，剩余关系
$\mathrm d(\partial_i f)=\sum_j(\partial_j\partial_i f)\,\mathrm dt_j$
给出$\Omega^1_{V^1/Y}$的一个方阵表示，矩阵是 Hessian（是否取转置取决于矩阵约定）。因此，若用
$\operatorname{Fitt}_0(\Omega^1_{V^1/Y})$
定义源端分歧子概形，其局部理想是$\det(\partial_i\partial_j f)$生成的主理想。

另一方面，二次初始项的极化在相同坐标中的矩阵也是该 Hessian：平方单项式$cT_i^2$在对角位置贡献$2c$，混合单项式$cT_iT_j$在对应两个非对角位置各贡献$c$。这里没有除以$2$，也没有把特征$2$的极化非退化改作二次零超曲面光滑。故该理想与$D$的零理想局部相同；Fitting 理想和线丛截面的零理想均可粘合，遂在这个明确的分歧结构下得到所拟的概形等式。
\endgroup

\endgroup

